% Abstract — ReCAHS paper
% Every number is traceable to cluster_results/runs/analysis/.

\begin{abstract}
Attention-head pruning methods almost universally score each head with a
single global importance value, and input-adaptive alternatives are
typically compared against static baselines that differ in \emph{both} the
selection criterion and the adaptivity --- leaving the source of any gain
unidentified. We present a controlled study of regime-conditional head
selection for time-series Transformers, using STL-derived
trend/seasonal/residual regimes on PatchTST across 4~datasets,
4~horizons, 5~seeds and 6~sparsity levels. Three findings emerge. First,
the standard global importance ranking is fragile: at realistic sparsity
it degrades sharply and is beaten by \emph{random} pruning on
minute-resolution data, whereas greedy backward elimination remains within
${\sim}1\%$ of the unpruned model at $25\%$ sparsity --- and with the mask
imposed from scratch during training, this pruning becomes free. Second,
regime-\emph{routing} adds nothing once the criterion is controlled: a
pooled greedy mask matches or beats regime-routed masks, and although a
per-window oracle over the regime masks outperforms even the unpruned
model in 15/16 cells, twelve deployable routers fail to capture this
headroom. Near-oracle in-sample fits, flat learning curves and failed
zero-shift diagnostics show the per-window optimal mask is aleatoric at
inference time. Third, the headroom is partially recoverable without
selection: \emph{averaging} the specialist masks' predictions beats the
best static mask on all four datasets and surpasses the unpruned model on
two --- an effect driven by mask diversity, for which regime partitioning
is the strongest, but not a necessary, source. On iTransformer's
cross-variate attention the criterion and ensemble findings replicate,
head redundancy is far larger, and --- unlike temporal attention --- a
small regime-routing benefit survives the control, marking a precise
architecture boundary. We release per-window losses for every
configuration to support further routing research.
\end{abstract}
